\section{Datenbasis und Datenaufbereitung}
\label{sec:datenbasis_und_datenaufbereitung}

Die Datenbasis und Datenaufbereitung sind entscheidende Schritte bei der Entwicklung eines Modells zur Vorhersage der Aluminiumpreise. In diesem Kapitel wird die Herkunft der Daten, die Datenformate, die Datenqualität und typische Fehler sowie die Charakteristika von Preisdaten näher erläutert.

Für die Analyse werden besteht die Datenbasis aus historischen Rohstoffpreisen und einem Aktienindex, der als Indikator für die wirtschaftliche Entwicklung dient. Betrachtet werden die Metalle Aluminium, Kupfer, Gold, Blei, Nickel, Silber und Zink sowie der Aktienindex S\&P 500. Die Metallpreise stammen aus der Datenbank von Naveen Sharma \cite{naveennas_metal_price_prediction_2024}, öffentlich zugänglich auf Kaggle, und umfassen tägliche Preise von Januar 2014 bis August 2024. Die Daten sind im CSV-Format vorliegend, was eine einfache Handhabung und Analyse ermöglicht. Sie enthalten Informationen wie das Datum, den Eröffnungspreis, den Höchstpreis, den Tiefstpreis, den Schlusskurs und das Handelsvolumen für jedes Metall. 

*asdfasdfdsafsadfd*





\begin{itemize}
    \item Umfang: ca. 1,5--2 Seiten
    \item Herkunft der Daten:
    \begin{itemize}
        \item synthetisch
        \item generiert
        \item offene Daten
    \end{itemize}
    \item Datenformate:
    \begin{itemize}
        \item CSV
        \item Zeitreihen
    \end{itemize}
    \item Datenqualität und typische Fehler
    \item Charakteristika von Preisdaten:
    \begin{itemize}
        \item Trends
        \item Saisonalität
        \item Ausreißer
        \item Strukturbrüche
    \end{itemize}
\end{itemize}